\section{Introduction}

\lai{I think we need a catchier intro paragraph (and especially opening sentence), but I'm not sure if the one below should be reworked or broken in half, with the FLP explanation in the second paragraph and the chatter about big ideas and contributions in the first. Something to consider.}

In functional logic programming, one program contains two semantics: a
``forward'' nondeterministic semantics with \emph{choices}, and a ``backward''
demand-driven semantics with \emph{free variables}. Past research has shown the
usefulness of this feature in different contexts.\ly{TODO: Citations.} In this
paper, we demonstrate yet another useful application: lazy cost analysis.
Specifically, we show that two different semantics proposed independently for
reasoning about the computation cost of lazy functional programs, 
clairvoyance semantics~\cite{clairvoyant,forking-paths} and demand
semantics~\cite{bjerner1989,demand-semantics}, can both be represented using a single functional logic program.

% \subsection{what is lazy evaluation?}

In lazy evaluation, a computation is driven by its demands, performing only the operations strictly required to compute a final result. Computations are
stored as \emph{thunks} which are only evaluated when their values are demanded, and only to the extent required to produce the demanded value.
This prevents the execution of extraneous code, but results in infamously challenging cost analysis. For example\lai{Gotta start with take n xs, where xs is arbitrarily long}, suppose that we define \emph{take n xs} as a function that takes the first n prefixes from a list xs.\lai{Should we give an actual definition, or leave it abstract like this to save space?} Let \([1..]\) represent an arbitrarily long list of natural numbers indexed at 1. Then, suppose that we wished to compute the following:

\lai{adding figures later because I'm just trying to get words down rn - for now, just text will do.}

\(take n [1..])\)

At first glance, one unfamiliar with laziness\lai{I think it's funnier to say this than the arguably more correct "one unfamiliar with lazy evaluation"} might expect this to fail to terminate, since defining [1..] would fail to terminate. However, since only n elements are demanded, so long as n is not arbitrarily long, the program will terminate in \(O(n)\) time. This is because n bounds the number of elements of [1..] that need to be computed, preventing the rest of the list from being generated. Thus, lazy programs can sometimes achieve termination when eager ones cannot.

\lai{This paragraph may be extraneous, but let's see how we are on space before ripping it out. It's good to have some mention of the difficulties here.}

One major drawback of laziness becomes apparent when we consider how partially executing a function call interferes with traditional cost analysis. Shortcutting (or even failing to run) functions runs counter to the concept of \(O\)-notation in particular. Technically speaking, every function takes \(O(1)\) time in a lazy language, as it is always possible to call a function and then never demand its output. For example, if we call \(take n [1..])\) as part of a larger function that never uses the resulting list of numbers, then the call would be stored in a \emph{thunk} and not further evaluated. As \emph{thunk} creation is a constant-time operation, this results in the aforementioned \(O(1)\) time cost. This is not a particularly \emph{meaningful} thing to say about a lazy function, but its truth renders much of the existing paradigm for cost analysis nearly useless for our purposes.

This paper will focus on two cost models which address this problem - \emph{Demand Semantics} and \emph{Clairvoyance Semantics}. Some other models will be explored briefly in Section whatever-number-is-related-work\lai{TODO: What number is this, anyway?}, but are not the focus of our work.

Clairvoyance Semantics models a lazy program as a nondeterministic tree of possible execution paths. The concept was first introduced in \cite{clairvoyant} and implemented in Rocq in \cite{forking-paths}\lai{Is this how we will do inline citations, and how we will reference them in the text?} via the Clairvoyance Monad. To analyze cost within this model, we utilize this monad to simulate nondeterministically choosing which thunks to force and which to skip, eventually resulting in a path which mimics that of a function computing the desired output. \lai{TODO: add more detail and corrections here - maybe we can use figures like those in the Forking Paths paper?} \lai{Do we need a concrete example here? Or do we need to mention that computing the actual tree is impractical because of path explosion?}

Demand Semantics, on the other hand, models the cost lazy programs by computing how much of the input must be forced to produce a function's output, which is provably cost-equivalent to computing the same output. First described in \cite{bjerner1989} and formalized in Rocq in \cite{demand-semantics}, the idea is also relatively straightforward, but is quite different from clairvoyance semantics in implementation. Rather than involving nondeterminism directly, the 'choices' made in a function's execution path are abstracted away into an "outD" input, which models the output which will ultimately be demanded from that function. outD is then fed to a separate cost-analysis function derived from the original function in a systematic way. Finally, the cost analysis function uses outD to compute the answer to a new query - how much of the original function's input would be needed to compute outD? It turns out that this computation results in the same time cost as computing the input normally.\lai{How about here?}

While clairvoyance and demand semantics provably produce equivalent cost analyses, directly translating from one to the other is challenging.\lai{Maybe two paragraphs - one outlining the problem, and another introducing the idea of using FLP to solve it.} There is certainly some amount of 'clairvoyance' occurring in demand semantics in the literal sense, as one needs to know what demands will be placed on the output before cost analysis can proceed in order to construct a meaningful outD. However, the handling of nondeterminism is left implicit - we take as premise that one will know which sorts of output demands a user might wish to reason about, and say little about how they came to possess that information. Clairvoyance semantics has a similar issue, but in reverse\footnote{After we introduce the main ideas of the functional logic translation between the two, the fact that clairvoyance and demand semantics are inverses of each other in this way will likely seem somewhere between ironic and prophetic.}. The concept of "demand" is not explicitly mentioned in the semantics, but nondeterministically choosing a program's execution path is also, \emph{implicitly}, very similar to computing that same path using the output demand as a guide \lai{Feels like I'm not explaining this quite well enough. I'll think about it.}.



While this conundrum left us (as researchers interested in lazy cost analysis) at an impasse\lai{too informal?}, those familiar with functional logic programming concepts may have already noticed the connection that eventually led us through it. Functional logic languages which integrate nondeterminism, like Curry\lai{What do I cite here...?}, have a notion of forward and backward execution that maps neatly onto the aforementioned relationship between demand and clairvoyance semantics. Is this mapping correct, though? Yes, we have discovered, and provably so.

We make the following contributions:

\begin{itemize}
  \item We propose and formalize a mechanical way of translating a functional
        program to a functional logic program to obtain the clairvoyance
        semantics and the demand semantics~(\cref{sec:translation}).
  \item We formally prove that such translation is correct with respect to both
        the clairvoyance semantics and the demand semnatics. All the theorems
        and proofs are mechanized in Agda~(\cref{sec:equiv}).
  \item We conduct a number of case studies by manually translating classical
        lazy Haskell programs\ly{TODO:\ replace these with exact names when we
        have the case study section.} to Curry/PEKCS programs based on our
        translation strategy~(\cref{sec:case-study}).
\end{itemize}

In addition, we provide a motivating example in \cref{sec:motivating-example},
discuss related work in \cref{sec:related-work}, and conclude this paper in
\cref{sec:conclusion}.
